%----------------------------------------------------------------------------------
% Exemplo do uso da classe tcc.cls. Veja o arquivo .cls
% para mais detalhes e instruções.
%----------------------------------------------------------------------------------
\chapter{\label{chap:related}Trabalhos Relacionados}

Este capítulo apresenta trabalhos relacionados ao uso de \textit{fuzzing} na verificação funcional de processadores RISC-V. Os trabalhos foram selecionados por abordarem diferentes aspectos necessários à construção de uma campanha de teste: a exploração de estados internos do RTL, a geração de programas válidos e complexos, a diversificação dos operandos e dos resultados das instruções, a mutação de programas em nível de blocos e a avaliação dos próprios \textit{fuzzers} e dos critérios que os orientam. 
A apresentação segue a ordem cronológica das publicações:
\begin{itemize}
    \item A \cref{sec:difuzz} apresenta o DifuzzRTL;
    \item a \cref{sec:cascade} apresenta o Cascade;
    \item a \cref{sec:encarsia} apresenta o Encarsia;
    \item a \cref{sec:divefuzz} apresenta o DiveFuzz;
    \item a \cref{sec:simfuzz} apresenta o SimFuzz;

\end{itemize}
% iniciando pelo DifuzzRTL e prosseguindo com Cascade, DiveFuzz e SimFuzz. 
Finalmente, a \cref{sec:compare} compara as abordagens com o trabalho proposto.

\section{DifuzzRTL}\label{sec:difuzz}

O DifuzzRTL \citep{difuzzrtl} é uma infraestrutura de \textit{fuzzing} diferencial que combina duas tarefas: conduzir implementações RTL a estados ainda não explorados e verificar se o comportamento observado está de acordo com a ISA. Suas entradas reúnem instruções, dados e interrupções. As instruções são geradas por mutação estruturada, com seleção aleatória de operações e operandos e reutilização de registradores e valores para favorecer dependências de dados.

A exploração é orientada pela cobertura dos registradores que influenciam sinais de controle do circuito. Durante a simulação, os valores desses registradores são combinados e registrados a cada ciclo; entradas que produzem combinações inéditas são preservadas no corpus. A correção, por sua vez, é verificada após a execução pela comparação dos registradores e da memória do RTL com o estado produzido por um simulador da ISA. A abordagem encontrou 16 falhas até então desconhecidas.

Para este trabalho, o DifuzzRTL evidencia a separação entre o critério que orienta a exploração e o oráculo que verifica a correção. A cobertura pode determinar quais entradas devem ser priorizadas, mas a identificação de divergências permanece sob responsabilidade da execução diferencial. Como limitação, a adoção dessa métrica exigiria uma análise e uma instrumentação específicas do RTL, não previstas no fluxo atual de validação do RS5.

\section{Cascade}\label{sec:cascade}

O Cascade \citep{cascade} produz programas RISC-V longos, válidos e com fluxos de dados e de controle interdependentes. A abordagem procura evitar que desvios não controlados encerrem prematuramente a execução e reduzam a quantidade de instruções efetivamente destinadas ao teste. Em vez de modificar um corpus, o Cascade gera sequências de blocos básicos e favorece o uso de registradores produzidos recentemente, formando cadeias de dependência.

O mecanismo central é uma pré-simulação da ISA que permite conhecer valores utilizados por saltos, acessos à memória e outras operações dependentes dos operandos. A partir dessas informações, o programa final é ajustado para preservar o fluxo esperado e, ao mesmo tempo, fazer com que decisões de controle dependam dos resultados anteriores. Uma divergência pode, assim, propagar-se até o fluxo de controle e impedir a terminação. O trabalho encontrou 37 novas falhas em processadores e uma falha no sintetizador Yosys.

O Cascade demonstra a importância de construir programas cujas instruções sejam efetivamente executadas e de formar dependências intencionais entre operações. Esses princípios podem ser aproveitados sem reproduzir integralmente a pré-simulação. Entretanto, a não terminação não identifica divergências que não alcancem o fluxo de controle; por isso, a comparação de assinaturas fornecida pelo RISCOF permanece necessária neste trabalho.

\section{Encarsia}\label{sec:encarsia}

O Encarsia \citep{encarsia} não é um \textit{fuzzer}, mas uma infraestrutura para avaliá-los. Os autores observam que cada \textit{fuzzer} é demonstrado em processadores diferentes e otimiza uma métrica própria, e que o número de falhas novas encontradas é um critério frágil de comparação, pois as falhas se tornam mais raras à medida que os projetos amadurecem. A partir de um levantamento de 177 correções em quatro processadores RISC-V de código aberto (Ibex, CVA6, Rocket e BOOM), o trabalho constata que as falhas observáveis pelo software se reduzem a dois padrões: trocas de sinais em atribuições e expressões e erros em comandos condicionais.

Com base nessa observação, o Encarsia injeta falhas automaticamente na representação intermediária do Yosys, trocando conexões entre sinais ou removendo e alterando casos das árvores de multiplexadores que representam as atribuições condicionais. Um circuito que compara o projeto original ao modificado é submetido à verificação formal, de modo que apenas falhas com efeito observável nos registradores arquiteturais entram no corpus. Com 90 falhas assim obtidas em Ibex, Rocket e BOOM, os autores comparam DifuzzRTL, ProcessorFuzz e Cascade sob as mesmas condições. Nenhum dos três detectou mais de 42\% das falhas em 24 horas; DifuzzRTL e ProcessorFuzz encontraram exatamente o mesmo conjunto de falhas com e sem a orientação por cobertura; e os três detectaram conjuntos distintos de falhas, diferença que os autores atribuem aos programas gerados, e não ao mecanismo de realimentação.

Para este trabalho, o Encarsia é relevante em dois pontos. Primeiro, propõe que a medida adequada de um \textit{fuzzer} é a detecção de falhas conhecidas, e não a cobertura alcançada. Segundo, mostra que métricas estruturais, como os registradores de controle do DifuzzRTL e as transições de CSR do ProcessorFuzz, não alteraram as falhas encontradas, e aponta como direção futura a realimentação funcional a partir de componentes da microarquitetura. Essas conclusões orientaram a escolha das fontes de sinal consideradas para o \textit{fuzzer} do RS5, discutida no \cref{chap:bg}.

\section{DiveFuzz}\label{sec:divefuzz}

O DiveFuzz \citep{divefuzz} busca aumentar a diversidade dos estados arquiteturais exercitados durante o \textit{fuzzing}. O trabalho observa que programas distintos podem produzir repetidamente os mesmos valores e que a frequência das instruções efetivamente executadas pode permanecer desequilibrada, mesmo quando os códigos de operação são gerados com probabilidades semelhantes.

Para reduzir essa repetição, os operandos são escolhidos durante a interação com um simulador da ISA, considerando o estado corrente dos registradores e da memória. O gerador favorece valores ainda pouco explorados e relações de dependência entre instruções, além de ajustar as probabilidades de geração para priorizar instruções raras. A correção é verificada pela comparação dos resultados produzidos pelo processador com o simulador. A abordagem aumentou a diversidade desses resultados e identificou 26 novas falhas.

O DiveFuzz mostra que a diversidade não deve ser medida apenas pelos códigos de operação, mas também pelos operandos, resultados e relações de dependência. Essa observação é relevante para o RS5, no qual diferentes relações entre instruções podem exercitar mecanismos distintos do pipeline. Contudo, a estratégia de comparação do DiveFuzz exigiria uma interface de co-simulação que não faz parte do fluxo atual do RISCOF; neste trabalho, o oráculo continuará baseado em assinaturas ao final da execução.

\section{SimFuzz}\label{sec:simfuzz}

O SimFuzz \citep{SIMFUZZ} explora o espaço de programas por meio de mutações em blocos orientadas por similaridade. A abordagem busca preservar a estrutura das sementes e reduzir a dependência de métricas de cobertura. Seu corpus inicial é composto por casos históricos que já revelaram falhas, sob a premissa de que esses programas representam regiões relevantes do espaço de entradas.

Cada semente é dividida em blocos delimitados por transferências de controle. As instruções internas podem ser substituídas, mas as saídas dos blocos são preservadas para manter o fluxo original. As substituições são limitadas por uma medida de similaridade que considera formato, código de operação, unidade funcional e operandos. Os programas mutados são comparados com um simulador da ISA, enquanto a cobertura é utilizada somente na avaliação. O SimFuzz identificou 17 falhas, das quais 14 eram desconhecidas.

O SimFuzz relaciona-se diretamente com este trabalho por adotar mutações em nível de bloco, preservar o fluxo de controle e utilizar a cobertura como informação complementar. A diferença principal está na origem do corpus: o SimFuzz parte de casos históricos de falhas, enquanto este trabalho utiliza blocos extraídos de aplicações reais e busca preservar seus padrões computacionais e relações de dependência.

\section{\label{sec:compare}Comparação com os trabalhos relacionados}

A Tabela~\ref{tab:comparacao-relacionados} sintetiza as características dos cinco trabalhos e as contrapõe à abordagem proposta. A estratégia de geração ou mutação indica como o espaço de programas é explorado, enquanto o critério de seleção ou feedback determina quais entradas são priorizadas a partir das informações disponíveis. O oráculo é apresentado separadamente, pois somente ele determina se o comportamento observado diverge da ISA.

\begin{table}[H]
    \centering
    \caption{Comparação dos trabalhos relacionados.}
    \label{tab:comparacao-relacionados}
    \small
    \setlength{\tabcolsep}{3pt}
    \resizebox{\textwidth}{!}{%
    \begin{tabular}{|p{1.8cm}|p{2.8cm}|p{3.2cm}|p{3.2cm}|p{3.2cm}|}
        \hline
        \textbf{Trabalho} & \textbf{Processadores} & \textbf{Estratégia de geração ou mutação} & \textbf{Critério de seleção ou feedback} & \textbf{Oráculo} \\
        \hline
        DifuzzRTL & Mor1kx e Rocket em ordem; BOOM superescalar fora de ordem & Geração e mutação individual de instruções e interrupções & Cobertura sensível a ciclos dos registradores de controle & Estado final do RTL comparado ao simulador da ISA \\
        \hline
        Cascade & PicoRV32, Kronos, VexRiscv, Rocket e CVA6 em ordem; BOOM superescalar fora de ordem & Geração de programas sintéticos longos em blocos básicos, com pré-simulação & Não utiliza feedback de execução; a pré-simulação orienta a construção & Não terminação causada pela propagação de divergências ao fluxo de controle \\
        \hline
        Encarsia & Ibex e Rocket em ordem; BOOM superescalar fora de ordem, com falhas injetadas & Não gera programas; injeta falhas no RTL (trocas de sinais e condicionais quebradas) para avaliar DifuzzRTL, ProcessorFuzz e Cascade & Avalia os critérios dos demais; a cobertura estrutural não alterou as falhas detectadas & Execução diferencial dos \textit{fuzzers} avaliados, com uma falha injetada por projeto \\
        \hline
        DiveFuzz & Rocket, CVA6 e NutShell em ordem; XiangShan superescalar fora de ordem & Construção dinâmica de instruções com operandos dependentes do contexto & Diversidade dos resultados e balanceamento dos códigos de operação & Comparação dos resultados com um simulador da ISA \\
        \hline
        SimFuzz & Rocket em ordem; BOOM e XiangShan superescalares fora de ordem & Mutação em blocos de casos históricos que revelaram falhas & Similaridade entre instruções; cobertura utilizada apenas na avaliação & Comparação diferencial com simuladores da ISA \\
        \hline
        Este trabalho & RS5 com pipeline em ordem, sob diferentes configurações & Combinação e mutação estruturada de blocos extraídos de aplicações reais & Diversidade estrutural e cobertura do RTL como critério complementar & Assinaturas do RS5 e do Sail comparadas pelo RISCOF \\
        \hline
    \end{tabular}
    }
\end{table}

A comparação evidencia duas estratégias predominantes. DifuzzRTL e DiveFuzz geram instruções em granularidade fina e utilizam informações da execução para direcionar novas entradas. Cascade e SimFuzz preservam uma estrutura maior dos programas, mas partem, respectivamente, de sequências sintéticas e de casos históricos de falhas. O Encarsia ocupa um papel distinto: não gera entradas, mas fornece o meio de comparar as demais abordagens e questiona o valor dos critérios estruturais que as orientam. As abordagens também distinguem os critérios que orientam a exploração dos mecanismos empregados para verificar a correção: cobertura, diversidade ou similaridade orientam a construção ou a escolha das entradas, enquanto a execução diferencial ou a não terminação identifica possíveis falhas.

Este trabalho combina elementos dessas estratégias em uma infraestrutura voltada ao RS5. Assim como no SimFuzz, as mutações ocorrem em nível de bloco e procuram preservar o fluxo de controle. A formação de dependências e a variação de operandos são motivadas pelo Cascade e pelo DiveFuzz. A diferença está no corpus composto por blocos de aplicações reais, em vez de programas inteiramente sintéticos ou casos históricos, e no uso conjunto da diversidade estrutural e da cobertura para selecionar entradas. O oráculo não constitui uma nova contribuição: a correção é determinada pela comparação de assinaturas entre o RS5 e o Sail já fornecida pelo RISCOF.

Outra lacuna observada é que os trabalhos comparados não têm como foco a avaliação sistemática de diferentes configurações de uma mesma implementação. A campanha proposta executará os programas gerados em combinações de extensões da ISA, políticas de encaminhamento, tamanhos de fila de instruções, organizações de memória e latências de resposta. Desse modo, a avaliação buscará verificar não apenas a conformidade das instruções, mas também sua interação com os diferentes caminhos internos disponibilizados pela parametrização do RS5.
